\chapter{Conclusion}
\label{ch:sum}

This thesis presented the design, implementation, and integration of a code coverage visualization feature for the CodeChecker static analysis platform.

\section{Summary of Contributions}

The implementation encompasses several key components:

\begin{itemize}
    \item \textbf{LCOV report converter}: A new \texttt{report-converter} module (\texttt{-t lcov}) that parses LCOV \texttt{.info} files and extracts covered lines, uncovered lines, and function coverage statistics into CodeChecker's \texttt{coverage.json} format. The converter follows the same auto-discovery pattern as all other analyzer integrations.

    \item \textbf{Database schema}: Two tables---\texttt{test\_coverage} for line-level data (using sign encoding to distinguish covered from uncovered lines) and \texttt{test\_coverage\_summary} for per-file function coverage statistics---with proper foreign key relationships and indexing.

    \item \textbf{API endpoints}: Two Thrift API endpoints (\texttt{getTestCoverage} and \texttt{getFilesWithCoverage}) that provide coverage data to the frontend. The \texttt{getFilesWithCoverage} endpoint is optimized to return line counts and function statistics in a single SQL query, avoiding per-file API calls.

    \item \textbf{Web interface}: Two Vue.js components---a directory-level statistics dashboard with LCOV-style coverage bars, function coverage columns, and file search; and a file-level view with line-by-line green/red/white highlighting based directly on LCOV instrumentation data.

    \item \textbf{Architectural decision}: The removal of the \texttt{CodeChecker coverage} CLI command in favor of the \texttt{report-converter} approach, aligning with CodeChecker's design as a static analysis tool and following the established pattern for external tool integration.
\end{itemize}

\section{Technical Achievements}

\begin{itemize}
    \item \textbf{Accurate line coloring}: By storing all LCOV \texttt{DA:} records (both covered and uncovered) and using the compiler's instrumentation as the authoritative source, the visualization matches \texttt{gcov}/\texttt{lcov} HTML reports exactly, eliminating the inaccuracies of heuristic-based executable line detection.

    \item \textbf{Scalable statistics}: The single-query optimization for the statistics page handles projects with hundreds of files (tested with 580 files from Apache Xerces-C++) without browser resource exhaustion.

    \item \textbf{Function coverage}: Extraction and display of function-level coverage statistics (\texttt{FNF}/\texttt{FNH}) from LCOV data, providing an additional dimension of coverage information beyond line counts.

    \item \textbf{Robust file matching}: Six fallback strategies for matching coverage file paths to database file IDs, handling differences in absolute paths, path prefixes, and filename-only matches.
\end{itemize}

\section{Future Work}

\begin{itemize}
    \item \textbf{Branch coverage}: Extending the feature to support branch coverage metrics (LCOV \texttt{BRDA:} records) would provide more comprehensive coverage information.

    \item \textbf{Coverage history}: Tracking coverage trends over time across code revisions would enable monitoring of test suite evolution.

    \item \textbf{Deeper static analysis integration}: Highlighting code regions with both static analysis warnings and low coverage would help prioritize areas for improvement.

    \item \textbf{Multiple coverage formats}: Supporting additional formats beyond LCOV (such as LLVM's \texttt{llvm-cov} JSON export) would increase applicability.

    \item \textbf{Execution counts}: Storing and displaying actual execution counts (not just covered/uncovered) would enable identification of hot paths and rarely-tested code.

    \item \textbf{Alembic migration for TestCoverageSummary}: Creating a proper database migration script for the function coverage summary table.
\end{itemize}

\section{Conclusion}

The code coverage visualization feature successfully extends CodeChecker by integrating test coverage metrics with static analysis results. The \texttt{report-converter} approach respects CodeChecker's design as a static analysis tool while providing a clean integration point for dynamic analysis data. The implementation was validated on a real-world C++ project (Apache Xerces-C++, 580 source files, 80 test cases) and produces accurate, performant results that match industry-standard coverage tools.
